\documentclass[english, parskip=full]{scrartcl}
\usepackage[utf8]{inputenc}  % Kodierung der Datei
\usepackage[ngerman]{babel}  % Sprache auf Deutsch 
\usepackage{color}
\usepackage{graphicx, gensymb}
\usepackage{amsfonts, cancel, amssymb}
\usepackage{amsmath, amsthm, array, esint}
\usepackage{booktabs, multirow}  % schönere Tabellen
\usepackage{caption}  % Anpassung der Captions
\usepackage{scrlayer-scrpage}    % Manipulation des Seitenstils
%\pagestyle{scrheadings}  % Stil für die Seite setzen
%\automark{section}  % Abschnittsnamen als Seitenbeschriftung 
%\setheadsepline{.5pt}    % Kopzeile durch Linie abtrennen
\usepackage[hidelinks]{hyperref}  % Links und weitere PDF-Features
\usepackage{typearea, tikz, pgffor, verbatim}
\usetikzlibrary{calc, quotes}
\usepackage[cache=false, outputdir=build]{minted}
\usepackage{placeins} % provides \FloatBarrier

\definecolor{bg}{rgb}{0.9,0.9,0.9}
\newcommand\numberthis{\addtocounter{equation}{1}\tag{\theequation}}
\usepackage{siunitx}
\newcommand{\bra}{}
\newcommand{\kom}{}
\newcommand{\brak}{}
\newcommand{\braket}{}
\newcommand{\intx}{}
\newcommand{\intr}{}
\newcommand{\kla}{}
\newcommand{\klam}{}
\newcommand{\klammer}{}
\newcommand{\oper}{}
\newcommand{\tecxt}{}
\newcommand{\Bra}{}
\newcommand{\Ket}{}
\renewcommand{\i}{\text{i}}
\newcommand{\e}{\text{e}}
\newcommand*{\Fourier}{\textsc{Fourier}\ }
\newcommand*{\F}{\mathcal{F}}
\newcommand*{\sinc}{\text{sinc\,}}
\newcommand{\conv}{\mathop{\scalebox{3.5}{\raisebox{-0.38ex}{\makebox[0.5\width][c]{$\ast$}}}}\limits}

\newcommand*{\titel}{04 Eigenwerte und Vektoriteration}
\newcommand*{\autor}{Joachim Schwardt und Julian Fleck}

\hypersetup{pdfauthor={\autor}, pdftitle={\titel}}

%\titlehead{titlehead}
\subject{Numerik II - Aufgaben}
\title{\titel}
%\author{\autor}
\author{\begin{tabular}{l|l}
Joachim Schwardt & 4768711 \\ 
Julian Fleck & 4759587\\ 
\end{tabular}}
\date{\today}

\begin{document}

%\begin{titlepage}
\maketitle  % Titel setzen
%\tableofcontents  % Inhaltsverzeichnis setzen
%\end{titlepage}

\section*{Aufgabe 1}
Gegeben ist eine Tridiagonalmatrix $A\in\mathbb{R}^{n\times n}$ mit
\begin{align}
A_{ij} &= \begin{cases} a & \tecxt\text{für } i=j+1,\\
b & \tecxt\text{für } i=j,\\
c & \tecxt\text{für } i=j-1,\\
0 & \tecxt\text{sonst}, \end{cases}
\end{align}
wobei $a,b,c\in\mathbb{R}$ und $ac > 0$.
Die Eigenwerte seien $\lambda_k$ und die Eigenvektoren seien $v_k$ für $k=1,\dots,n$.
Die Eigenwertgleichung lautet dann
\begin{align}
\begin{pmatrix}
b & c & 0 & \dots \\
a & b & c & \dots\\
0 & a & b & \dots\\
\vdots & \vdots & \vdots & \ddots
\end{pmatrix} \begin{pmatrix}
v_{k1} \\ v_{k2} \\ v_{k3} \\ \vdots
\end{pmatrix} &= \begin{pmatrix}
bv_{k1} + cv_{k2} \\ av_{k1} + bv_{k2} + cv_{k3} \\ av_{k2} + bv_{k3} + cv_{k4} \\ \vdots
\end{pmatrix} = \lambda_k \begin{pmatrix}
v_{k1} \\ v_{k2} \\ v_{k3} \\ \vdots
\end{pmatrix}
\end{align}
Die erste Gleichung ergibt $v_{k1} = \frac{c}{\lambda_k - b} v_{k2}$.
Einsetzen in die nächste Gleichung ergibt
\begin{align}
0 &= av_{k1} + (b-\lambda_k) v_{k2} + cv_{k3} = \frac{ac}{\lambda_k - b} v_{k2} - (\lambda_k - b) v_{k2} + cv_{k3},\\
\Rightarrow\ \ \ v_{k2} &= \frac{c(b-\lambda_k)}{ac - (\lambda_k - b)^2} v_{k3}
\end{align}
Mit $j=2,\dots,n-1$ lauten die nächsten Gleichungen dann
\begin{align}
av_{k,j-1} + bv_{k,j} + cv_{k,j+1} &= \lambda_k v_{k,j},\\
\Rightarrow\ \ \ v_{k,j+1} &= \frac{\lambda_k - b}{c} v_{k,j} - \frac{a}{c} v_{k,j-1}.
\end{align}
Hier ist es hilfreich zunächst allgemein die Iterationsvorschrift
\begin{align}
x_{j+1} &:= \beta x_j + \alpha x_{j-1}
\end{align}
zu betrachten.
Wir führen eine weitere Variable $y_{j} := x_{j-1}$ und definieren $X_j := (x_j,y_j)^\top$.
Dann lautet die Iteration
\begin{align}
X_{j+1} &= \underbrace{\begin{pmatrix}
\beta & \alpha \\ 1 & 0
\end{pmatrix}}_{=:\,C}  X_j
\end{align}
Das charakteristische Polynom ist $\chi_C(\mu) = \mu^2 - \beta\mu - \alpha$, die Eigenwerte lauten also $\mu_{1,2} = \frac{\beta}{2} \pm \sqrt{\frac{\beta^2}{4} + \alpha}$.
Wir finden die Eigenvektoren $u_{1,2}$ über
\begin{align}
\beta u_{k,1} + \alpha u_{k,2} &= \mu_k u_{k,1},\\
u_{k,1} &= \mu_k u_{k,2},
\end{align}
also $u_{k} = (\mu_k,1)^\top$.
Die Transformationsmatrix lautet dann
\begin{align}
T &:= \begin{pmatrix}
\mu_1 & \mu_2 \\ 1 & 1
\end{pmatrix},
\end{align}
wobei $C = TDT^{-1}$ mit $D=\tecxt\text{diag\,}(\mu_1,\mu_2)$.
Überprüfen:
\begin{align}
TDT^{-1} &= \begin{pmatrix}
\mu_1 & \mu_2 \\ 1 & 1
\end{pmatrix} \begin{pmatrix}
\mu_1 & 0 \\ 0 & \mu_2
\end{pmatrix} \frac{1}{\mu_1 - \mu_2} \begin{pmatrix}
1 & -\mu_2 \\ -1 & \mu_1
\end{pmatrix} = \\
&= \frac{1}{\mu_2 - \mu_1} \begin{pmatrix}
\mu_1 & \mu_2 \\ 1 & 1
\end{pmatrix} \begin{pmatrix}
\mu_1 & -\mu_1\mu_2 \\ -\mu_2 & \mu_1\mu_2
\end{pmatrix} = \\
&= \frac{1}{\mu_1 - \mu_2} \begin{pmatrix}
\mu_1^2 - \mu_2^2 & \mu_1\mu_2(\mu_2 - \mu_1) \\
\mu_1 - \mu_2 & 0
\end{pmatrix} = \begin{pmatrix}
\mu_1 + \mu_2 & -\mu_1\mu_2 \\ 1 & 0
\end{pmatrix} = C.
\end{align}
Damit können wir die Lösung explizit angeben.
Für $x_0=0$ ergibt sich
%\begin{align}
%X_{j} &= C^{j-2} X_2 = TD^jT^{-1} X_1 = T\begin{pmatrix}
%\mu_1^{j-2} & 0 \\ 0 & \mu_2^{j-2}
%\end{pmatrix} T^{-1} \begin{pmatrix}
%x_2 \\ x_1
%\end{pmatrix} = \\
%&= \begin{pmatrix}
%\mu_1 & \mu_2 \\ 1 & 1
%\end{pmatrix} \begin{pmatrix}
%\mu_1^{j-2} & 0 \\ 0 & \mu_2^{j-2}
%\end{pmatrix} \frac{1}{\mu_1 - \mu_2} \begin{pmatrix}
%1 & -\mu_2 \\ -1 & \mu_1
%\end{pmatrix} \begin{pmatrix}
%x_2 \\ x_1
%\end{pmatrix} = \\
%&= \frac{1}{\mu_1 - \mu_2} \begin{pmatrix}
%\mu_1 & \mu_2 \\ 1 & 1
%\end{pmatrix} \begin{pmatrix}
%\mu_1^{j-2} (x_2 - \mu_2 x_1) \\ -\mu_2^{j-2} (x_2 - \mu_1 x_1)
%\end{pmatrix} = \\
%&= \frac{1}{\mu_1 - \mu_2} \begin{pmatrix}
%(\mu_1^{j-1} - \mu_2^{j-1})x_2 - \mu_1^{j-1} \mu_2x_1 + \mu_2^{j-1} \mu_1x_1 \\
%(\mu_1^{j-2} - \mu_2^{j-2}) x_2 - \mu_1^{j-2}\mu_2 x_1 + \mu_2^{j-2}\mu_1 x_1)
%\end{pmatrix}
%\end{align}
\begin{align}
X_{j} &= C^{j-1} X_1 = TD^{j-1}T^{-1} X_1 = T\begin{pmatrix}
\mu_1^{j-1} & 0 \\ 0 & \mu_2^{j-1}
\end{pmatrix} T^{-1} \begin{pmatrix}
x_1 \\ x_0
\end{pmatrix} = \\
&= \begin{pmatrix}
\mu_1 & \mu_2 \\ 1 & 1
\end{pmatrix} \begin{pmatrix}
\mu_1^{j-1} & 0 \\ 0 & \mu_2^{j-1}
\end{pmatrix} \frac{1}{\mu_1 - \mu_2} \begin{pmatrix}
1 & -\mu_2 \\ -1 & \mu_1
\end{pmatrix} \begin{pmatrix}
x_1 \\ 0
\end{pmatrix} = \\
&= \frac{1}{\mu_1 - \mu_2} \begin{pmatrix}
\mu_1 & \mu_2 \\ 1 & 1
\end{pmatrix} \begin{pmatrix}
\mu_1^{j-1} x_1 \\ -\mu_2^{j-1} x_1
\end{pmatrix} = \\
&= \frac{1}{\mu_1 - \mu_2} \begin{pmatrix}
(\mu_1^j - \mu_2^j)x_1 \\ (\mu_1^{j-1} - \mu_2^{j-1})x_1
\end{pmatrix},
\end{align}
also insbesondere
\begin{align}
x_j &= \frac{\mu_1^j - \mu_2^j}{\mu_1 - \mu_2}x_1 = \sum_{k=0}^{j-1} \mu_2^k \mu_1^{j-k-1} x_1
\end{align}
Da die Eigenvektoren nur $n$ Einträge haben, fordern wir hier insbesondere noch $x_{n+1} = 0$.
Daraus folgt die Bedingung
\begin{align}
\mu_1^{n+1} &= \mu_2^{n+1}.
\end{align}
Da insbesondere $\mu_1 \neq \mu_2$, muss also $\mu_{1,2}\in\mathbb{C}$ gelten.
Wir schreiben daher $\mu_k = |\mu_k|\e^{\i\phi_k}$.
Dann ist die Bedingung äquivalent zu $|\mu_1| = |\mu_2|$ und $(n+1)\phi_1 \equiv (n+1)\phi_2\,\tecxt\text{mod}\,2\pi$.
Das ist nur möglich, wenn $\mu_{1,2} = R \pm \i I$ mit $R,I\in\mathbb{R}$ ist, also hier
\begin{align}
\mu_{1,2} &= \frac{\beta}{2} \pm \i\sqrt{-\alpha - \frac{\beta^2}{4}} \equiv R + \i I.
\end{align}
Betrag und Phase sind dann
\begin{align}
|\mu_{1,2}| &= \sqrt{-\alpha},\\
\phi_{1,2} &= \pm \arctan\kla\left( \frac{\sqrt{-\alpha - \frac{\beta^2}{4}}}{\frac{\beta}{2}} \right).
\end{align}
Eine mögliche Lösung der ``Phasengleichung'' wäre dann
\begin{align}
(n+1)\phi_{1,2} &= \pm k\pi,\\
\Rightarrow\ \ \ \frac{\sqrt{-\alpha - \frac{\beta^2}{4}}}{\frac{\beta}{2}} &= \tan\kla\left( \frac{k\pi}{n+1} \right),\\
\Rightarrow\ \ \ \beta &= 2\sqrt{-\alpha} \cos\kla\left( \frac{k\pi}{n+1} \right).
\end{align}
%Wir finden auf dem selben Weg auch eine einfachere Darstellung für die Phasen $\cos\phi_{1,2} = \frac{\beta}{2\sqrt{-\alpha}}$.
%\\
In unserem Fall ist $\beta = \frac{\lambda_k - b}{c}$ und $\alpha = -\frac{a}{c}$.
Die Eigenwerte sind also
\begin{align}
\lambda_k &= b + 2c\sqrt{\frac{a}{c}} \cos\kla\left( \frac{k\pi}{n+1} \right) = b + 2\tecxt\text{sgn\,}(c) \sqrt{ac} \cos\kla\left( \frac{k\pi}{n+1} \right),
\end{align}
wobei noch $\tecxt\text{sgn\,}(c) = \tecxt\text{sgn\,}(a)$ aus $ac > 0$ folgt.
Die Eigenwerte von $C$ lauten dann
\begin{align}
\mu_{1,2} &= \sqrt{\frac{a}{c}} \exp\kla\left( \i \phi_{1,2} \right) =  \sqrt{\frac{a}{c}} \exp\kla\left( \pm\i \frac{k\pi}{n+1} \right) .
\end{align}
Die Komponenten der Eigenvektoren ergeben sich dann zu
\begin{align}
v_{k,j} &= \frac{\mu_1^j - \mu_2^j}{\mu_1 - \mu_2} = \frac{\sqrt{\frac{a}{c}}^j}{\sqrt{\frac{a}{c}}} \frac{\e^{\i \frac{k\pi}{n+1} j} - \e^{-\i \frac{k\pi}{n+1} j}}{\e^{\i \frac{k\pi}{n+1} } - \e^{-\i \frac{k\pi}{n+1} }} x_1 = \\
&= \kla\left( \frac{a}{c} \right)^{\frac{j-1}{2}} \sin\kla\left( \frac{k\pi }{n+1}j \right) \frac{x_1}{\sin\kla\left( \frac{k\pi}{n+1} \right)} \equiv \kla\left( \frac{a}{c} \right)^{\frac{j-1}{2}} \sin\kla\left( \frac{k\pi}{n+1}j \right),
\end{align}
wobei wir $x_1 := \sin\kla\left( \frac{k\pi}{n+1} \right)$ definiert haben.
\\
Sei $b=2$ und $a=c=-1$.
Dann lauten die Eigenwerte
\begin{align}
\lambda_k &= 2 - 2\cos\kla\left( \frac{k\pi}{n+1} \right) = 4\sin^2\kla\left( \frac{k\pi}{2(n+1)} \right).
\end{align}
Die Kondition kann hier über das Verhältnis der Eigenwerte berechnet werden, also
%\begin{align}
%\kappa &= \frac{\max_k \lambda_k}{\min_k \lambda_k} = \frac{4\sin^2\kla\left( \frac{n\pi}{2(n+1)} \right)}{4\sin^2\kla\left( \frac{\pi}{2(n+1)} \right)} = \kla\left\vert \frac{\e^{\i \frac{n\pi}{2(n+1)}} - \e^{-\i \frac{n\pi}{2(n+1)}}}{\e^{\i \frac{\pi}{2(n+1)}} - \e^{-\i \frac{\pi}{2(n+1)}}} \right\vert^2 = \\
%&= \left\vert \frac{1 - \e^{-\i\frac{n\pi}{n+1}}}{1 - \e^{-\i\frac{\pi}{n+1}}} \right\vert^2 = \left\vert \sum_{k=0}^{n-1} \e^{-\i \frac{k\pi}{n+1}} \right\vert^2 = \\
%&= \left\vert \sum_{k=0}^{n} \e^{-\i \frac{k\pi}{n+1}} - \e^{-\i \frac{n\pi}{n+1}}\right\vert^2 = \left\vert \frac{2}{1 - \e^{-\i\frac{\pi}{n+1}}} - \e^{-\i \frac{n\pi}{n+1}} \right\vert^2 
%\end{align}
\begin{align}
\kappa &= \frac{\max_k \lambda_k}{\min_k \lambda_k} = \frac{4\sin^2\kla\left( \frac{n\pi}{2(n+1)} \right)}{4\sin^2\kla\left( \frac{\pi}{2(n+1)} \right)} = \\
&= \frac{\sin^2\kla\left( \frac{\pi}{2} - \frac{\pi}{2} \frac{1}{n+1} \right)}{\sin^2\kla\left( \frac{\pi}{2} \frac{1}{n+1} \right)} = \frac{1}{\tan^2\kla\left( \frac{\pi}{2} \frac{1}{n+1} \right)} \approx \frac{4(n+1)^2}{\pi^2}.
\end{align}
Für $n=9^2$ ist $\kappa \approx \SI{2.72e3}{}$ und für $n=99^2$ ist $\kappa \approx \SI{3.89e7}{}$.

\section*{Aufgabe 2}
Sei $A\in\mathbb{R}^{n\times n}$ mit den Eigenwerten $|\lambda_1|>|\lambda_2|\ge |\lambda_3|\ge \dots\ge |\lambda_n|$ gegeben.
Sei $v_1$ der Eigenvektor zum einfachen Eigenwert $\lambda_1$.
Außerdem sei $x^0\in\mathbb{R}$ mit $\brak\langle x^0 | v_1 \rangle \neq 0$.
Die Vektoriteration erzeugt Folgen $\kla\left( x^{(j)} \right)_{j\in\mathbb{N}}$ und $\kla\left( \lambda^{(j)} \right)_{j\in\mathbb{N}}$ mit
\begin{align}
x^{(j)} &:= A^j x^0, \quad \lambda^{(j)} &:= \frac{\braket\langle x^{(j)} | A | x^{(j)} \rangle}{\brak\langle x^{(j)} | x^{(j)} \rangle},
\end{align}
die gegen $v_1$ bzw. $\lambda_1$ konvergieren.
Sei $A$ zusätzlich symmetrisch.
Dann sind die Eigenvektoren orthonormal, also $\brak\langle v_j | v_k \rangle = \delta_{jk}$.
Wir entwickeln den Startvektor in dieser Basis mit den Koeffizienten $c_j := \brak\langle v_j | x^0 \rangle$, wobei nach Voraussetzung $c_1\neq 0$.
Dann gilt
\begin{align}
x^0 &= \sum_{k=1}^n c_kv_k,\\
x^{(j)} &= A^jx^0 = \sum_{k=1}^n c_k\lambda_k^j v_k.
\end{align}
Damit gilt für die Eigenwerte
\begin{align}
\lambda^{(j)} &= \frac{\braket\langle \sum_{k=1}^n c_k\lambda_k^jv_k | A | \sum_{l=1}^n c_l\lambda_l^jv_l \rangle}{\brak\langle \sum_{k=1}^n c_k\lambda_k^jv_k  | \sum_{l=1}^n c_l\lambda_l^{j}v_k \rangle} \kom\overset{\substack{{\brak\langle v_k | v_l \rangle\,=\,\delta_{kl}} \\ |}}{=} \frac{\sum_{k=1}^n |c_k|^2|\lambda_k|^{2j}\lambda_k}{\sum_{k=1}^n |c_k|^2|\lambda_k|^{2j}} = \\
&= \frac{\lambda_1 + \sum_{k=1}^n \left\vert \frac{c_k}{c_1}\right\vert^2 \left\vert \frac{\lambda_k}{\lambda_1} \right\vert^{2j}\lambda_k}{1 + \sum_{k=1}^n \left\vert \frac{c_k}{c_1} \right\vert^2 \left\vert \frac{\lambda_k}{\lambda_1} \right\vert^{2j}} \le \lambda_1 + \sum_{k=1}^n \left\vert \frac{c_k}{c_1}\right\vert^2 \left\vert \frac{\lambda_k}{\lambda_1} \right\vert^{2j}\lambda_k.
\end{align}
Damit finden wir
\begin{align}
|\lambda^{(j)} - \lambda_1| &\le \sum_{k=1}^n \left\vert \frac{c_k}{c_1}\right\vert^2 \left\vert \frac{\lambda_k}{\lambda_1} \right\vert^{2j}|\lambda_k| \in\mathcal{O}\kla\left( \left\vert \frac{\lambda_2}{\lambda_1} \right\vert^{2j} \right).
\end{align}
Im letzten Schritt nutzen wir aus, dass $|\lambda_2| \ge |\lambda_{k>2}|$, weshalb für $j\rightarrow\infty$ der Term mit $\left\vert \frac{\lambda_2}{\lambda_1} \right\vert$ gegenüber den anderen dominiert.

\section*{Aufgabe 3}
Gegeben ist das Anfanswertproblem $x'=f(x)$ mit $x(0)=1$.
Die Lipschitz-Konstante einer differenzierbaren Funktion auf einer Definitionsmenge $D$ ist gegeben durch
\begin{align}
L(f) &= \sup_{x\in D} |f'(x)|.
\end{align}
Wir untersuchen im Folgenden alle Funktionen auf Lipschitz-Stetigkeit in einer Umgebung von $x=1$, also $D= (1-\epsilon, 1+\epsilon)$ für ein $\epsilon>0$.
\subsection*{(a)} 
Für $f(x)=1-x^2$ gilt 
\begin{align}
L &= \sup_{x\in D} |-2x| = 2|1+\epsilon|,
\end{align}
die Funktion ist also Lipschitz-stetig.
\subsection*{(b)} 
Für $f(x)=|1-x^2| = \sqrt{(1-x^2)^2}$ gilt 
\begin{align}
L &= \sup_{x\in D} \left\vert \frac{1}{2\sqrt{(1-x^2)^2}} 2(1-x^2)(-2x) \right\vert = 2\sup_{x\in D}  \left\vert \frac{1-x^2}{|1-x^2|} x \right\vert = 2(1+\epsilon),
\end{align}
die Funktion ist also Lipschitz-stetig.
\subsection*{(c)} 
Für $f(x)= \sqrt{|1-x^2|} = [(1-x^2)^2]^{1/4}$ gilt 
\begin{align}
L &= \sup_{x\in D} \left\vert \frac{1}{4[(1-x^2)^2]^{3/4}} 2(1-x^2)(-2x) \right\vert = \sup_{x\in D}  \left\vert \frac{x}{\sqrt{|1-x^2|}} \right\vert = \infty,
\end{align}
die Funktion ist also bei $x=1$ nicht Lipschitz-stetig.
\\
Die Bedingung für Picard-Lindelöf ist gerade, dass $f$ Lipschitz-stetig im ``zweiten Argument'', hier also $x$ ist.
Damit kann man die eindeutige Lösbarkeit für (a) und (b) folgern.


\section*{Programmieraufgabe 4 (schriftlicher Teil)}

%\begin{thebibliography}{9}
%\bibitem{example} description, \url{https://www.exmapleurl}, day.\,month~2021.
%\end{thebibliography}

\end{document}
